\section*{Bandpass VCVS}

\noindent\colorbox{orange}{
  \begin{minipage}{1.0\linewidth}
    \paragraph{Richiesta (a)}
    Dimostrare, sulla base di semplici argomenti di natura qualitativa, che il circuito proposto
    svolga effettivamente la funzione di filtro passa-banda, discutendone la risposta ai limiti
    del dominio delle frequenze reali.
  \end{minipage}
}
\paragraph{Risposta (a)}
La tensione d'ingresso $V_S$ attraversa un arrangiamento di partitori configurati come
un filtro passa basso (passivo) e un filtro passa alto (passivo) collegati a cascata;
il filtro passa basso ha un ramo in comune con la retroazione dell'opamp.
L'opamp segue lo schema non invertente.

Al limite del dominio delle frequenze reali si manifestano le seguenti condizioni
\begin{itemize}
\item
  Per frequenze $f \to 0$ i condensatori si comportanto come rami aperti il nodo B
  \'e messo a terra e nel regime di operazione lineare
  \begin{equation*}
    V_{OUT} = A_d(0-V_-) = A_d(0-kV_{OUT})
  \end{equation*}
  la tensione all'ingresso dell'ingresso invertente ha un fattore di partizione positivo
  $k$ (non \'e necessario calcolarlo)
  \begin{equation*}
    A_d(0-kV_{OUT}) \implies V_{OUT} = \frac{1}{1+kA_d}
  \end{equation*}
  per guadagni differenziali elevati la tensione $\colorbox{blu}{\boxed{V_{OUT} \to 0}}$
\item
  Per frequenze $f \to \infty$ i condensatori si comportanto come corto-circuiti il nodo A
  \'e messo a terra, la tensione al nodo B \'e nulla quindi, per un argomento analogo al caso
  $f\to 0$, per guadagni differenziali elevati la tensione $\colorbox{blu}{\boxed{V_{OUT} \to 0}}$.
\end{itemize}

\noindent\colorbox{orange}{
  \begin{minipage}{1.0\linewidth}
    \paragraph{Richiesta (b)}
    Derivare analiticamente la funzione di trasferimento complessa Av(s)=VOUT(s)/VS(s)
    nell’ipotesi che l’operazionale sia ideale (\textit{suggerimento: si determini la tensione nel
      punto A del circuito in termini di VS e VOUT; indipendentemente, si esprimano VB e quindi
      VA in termini di Vout assumendo che l'opAmp funzioni in regime lineare ed utilizzando il
      principio di cortocircuito virtuale; si eguaglino infine le espressioni ottenute per ottenere
      la funzione di trasferimento}). Per semplicit\'a di calcolo, si assuma inoltre (ipotesi
    compatibile con i valori e le incertezze nominali dei componenti) che R1 = R3 = R, R2 =
    2 R, C1 = C2 = C, $R_B$ = 5/3 $R_A$.
  \end{minipage}
}
\paragraph{Risposta (b)}
La richiesta porta con se un suggerimento, seguendolo per il nodo A vale
\begin{equation*}
  \frac{V_S - V_A}{R} = sCV_A + sC(V_A - V_B) + \frac{V_A - V_{OUT}}{R}
\end{equation*}
quindi l'espressione della KCL per il nodo A include i termini $V_A$ e $V_{OUT}$
successivamente, sempre seguendo le istruzioni date si trovano ulteriori
relazioni per i nodi $V_A$ e $V_B$
\begin{equation*}
  sC(V_A - V_B) = \frac{V_B}{2R}
\end{equation*}
assumendo che l'opamp sia ideale e che operi in regime lineare si applica
il principio di corto-circuito virtuale $V_+ = V_-$ da cui
\begin{equation*}
  \begin{gathered}
    \begin{cases}
      \frac{V_{OUT} - V_-}{2R} = \frac{V_-}{\frac65R} \\
      V_B = V_+ = V_-
    \end{cases} \implies
    \frac{V_{OUT} - V_B}{2R} = \frac{V_B}{\frac65R} \\
    V_B = \frac38V_{OUT}
  \end{gathered}
\end{equation*}
ora $V_B$ \'e puramente espressione di $V_{OUT}$ questo permetter\'a
di trovare una relazione analoga per $V_A$
\begin{equation*}
  \begin{gathered}
    \begin{cases}
      sC(V_A - V_B) = \frac{V_B}{2R} \\
      V_B = \frac38V_{OUT}
    \end{cases} \implies
    sC\left(V_A - \frac38V_{OUT}\right) = \frac{\frac38V_{OUT}}{2R} \\
    V_A = \frac38\left(\frac{1 + 2RCs}{2RCs}\right)V_{OUT}
  \end{gathered}
\end{equation*}
ora $V_B$ e $V_A$ sono espressioni unicamente di $V_{OUT}$. Ora si riduce
la prima equazione nodale trovata all'inizio di questa risposta
\begin{equation*}
  V_S = (2 + 2sRC)V_A - sRCV_B -V_{OUT}
\end{equation*}
e si sostituiscono $V_A$ e $V_B$ con le loro espressioni in termini di $V_{OUT}$
\begin{equation*}
  V_S = \left[\frac38(2 + 2sRC)\left(\frac{1 + 2RCs}{2RCs}\right) - 1 - \frac38sRC\right]V_{OUT}
\end{equation*}
l'espressione finale per la funzione di trasferimento complessa
\begin{equation}
  \label{eq:bp-vcvs-tfunc}
  \colorbox{blu}{\boxed{H(s) = \frac{8RCs}{3 + RCs + 3R^2C^2s^2}}}
\end{equation}

\noindent\colorbox{orange}{
  \begin{minipage}{1.0\linewidth}
    \paragraph{Richiesta (c)}
    Dalla funzione di trasferimento dedurre i valori attesi per i parametri del filtro (frequenza
    centrale, guadagno massimo e larghezza banda) e confrontarli con le misure di cui ai
    punti 1a, 1b, 1c.
  \end{minipage}
}
\paragraph{Risposta (c)}
Le parti reali e immaginarie del numeratore e denominatore
\begin{equation*}
  \begin{cases}
    \mathcal{R}_{\omega}^N = 0 \\
    \mathcal{I}_{\omega}^N = 8RC \\
    \mathcal{R}_{\omega}^D = 3 - 3R^2C^2 \\
    \mathcal{I}_{\omega}^D = RC
  \end{cases} \implies
  G(\omega) = \sqrt{\frac{(8RC\omega)^2}{(3 - 3R^2C^2\omega^2)^2 + (RC\omega)^2}}
\end{equation*}
il guadagno \'e ottenuto dall'espressione \ref{eq:da-tfunc-a-gain}
\begin{equation}
  \label{eq:bp-vcvs-gain}
  G(f) = \frac{16\pi RCf}{\sqrt{9 - 68\pi^2R^2C^2f^2 + 144\pi^4R^4C^4f^4}}
\end{equation}
il guadagno massimo pu\'o essere ottenuto semplicemente riducendo l'espressione
del guadagno trovata in precedenza (in funzione della pulsazione $\omega$) nella
seguente forma
\begin{equation*}
  G(\omega) = \frac{8}{\sqrt{\left(\frac{3 - 3R^2C^2\omega^2}{RC\omega}\right)^2 + 1}}
\end{equation*}
il numeratore assume valore minimo per $3 - 3R^2C^2\omega^2 = 0$ e quindi il guadagno \'e
\begin{equation}
  \label{eq:bp-vcvs-gain-max}
  \colorbox{blu}{\boxed{\max_{\omega \in [0, +\infty)} G(\omega) = 8}}
\end{equation}
inoltre sostituendo $\omega \to 2\pi f$ alla condizione di minimizzazione si ricava la
frequenza centrale
\begin{equation}
  \label{eq:bp-vcvs-frequenza-centrale}
  \colorbox{blu}{\boxed{f = \frac{1}{2\pi RC}}}
\end{equation}
Un modo di ottenere la larghezza di banda richiede le frequenze di taglio, si calcolano
vincolando il guadagno quadro alla met\'a del suo guadagno massimo quadro (si usa il guadagno
in funzione della pulsazione per facilitare i calcoli)
\begin{equation*}
  \begin{gathered}
    \frac{8RC\omega}{\sqrt{9 - 17R^2C^2\omega^2 + 9R^4C^4\omega^4}} = \frac{8}{\sqrt2} \\ \implies
    9 - 19R^2C^2\omega^2 + 9R^4C^4\omega^4 = 0
  \end{gathered}
\end{equation*}
un ulteriore sostituzione $R^2C^2\omega^2 \to \xi$ semplifica il problema ulteriormente ad
un polinomio di ordine 2
\begin{equation*}
  \begin{gathered}
    9 - 19\xi + 9\xi^2 = 0 \\
    \xi = \frac{19 \pm \sqrt{37}}{18}
  \end{gathered}
\end{equation*}
sostituendo $\xi \to R^2C^2\omega^2$ e $\omega \to 2\pi f$ si ottiene
\begin{equation*}
  \begin{gathered}
    4\pi^2R^2C^2f^2 = \frac{19 \pm \sqrt{37}}{18} \implies
    f_{L,H} = \frac{1}{2\pi RC}\sqrt{\frac{19 \pm \sqrt{37}}{18}}
  \end{gathered}
\end{equation*}
dunque la larghezza di banda $\Delta f = f_H - f_L$
\begin{equation}
  \label{eq:bp-vcvs-bw}
  \colorbox{blu}{\boxed{\Delta f = \frac{1}{2\pi\sqrt{18}RC}\left[\sqrt{19 + \sqrt{37}} - \sqrt{19 - \sqrt{37}}\right]}}
\end{equation}

\noindent\colorbox{orange}{
  \begin{minipage}{1.0\linewidth}
    \paragraph{Richiesta (d)}
    Dare un’interpretazione qualitativa (e se possibile quantitativa) delle caratteristiche del
    segnale di uscita osservato al punto 1d. Che tipo di segnale appare? A quale ampiezza?
  \end{minipage}
}
